% !TeX program = xelatex
% 从工程根目录运行：latexmk examples/code-styles.tex
\documentclass[10pt,letterpaper]{book}
\input{config/preamble}
\newfontfamily\PreviousCodeFont{DejaVuSansMono.ttf}[
  Path=fonts/,Scale=0.9,Ligatures=NoCommon,
  BoldFont=DejaVuSansMono-Bold.ttf,
  ItalicFont=DejaVuSansMono-Oblique.ttf,
  BoldItalicFont=DejaVuSansMono-BoldOblique.ttf
]
\begin{document}
\pagestyle{plain}
\section*{代码字体与语法高亮示例}

下面两段使用相同代码和配色，可以直接比较字体。新版采用接近原稿的
Latin Modern Mono；中文注释使用 Fandol 仿宋风格。纯文本和目录树继续使用
DejaVu，以完整显示线条和特殊空格。

\subsection*{新版程序字体：Latin Modern Mono}
\begin{CodeBlock}[language=BookVerilog]
// 时序逻辑：低电平复位，保留 0/O、1/l、下划线和引号的区别
module counter(input wire clk, reset_n, output reg [7:0] count);
  always @(posedge clk) begin
    if (!reset_n) count <= 8'h00;
    else count <= count + 1'b1;
  end
  initial $display("counter is ready");
endmodule
\end{CodeBlock}

\subsection*{备选等宽字体：DejaVu Sans Mono（同配色对照）}
\begin{CodeBlock}[language=BookVerilog,basicstyle=\small\ttfamily\PreviousCodeFont]
// 时序逻辑：低电平复位，保留 0/O、1/l、下划线和引号的区别
module counter(input wire clk, reset_n, output reg [7:0] count);
  always @(posedge clk) begin
    if (!reset_n) count <= 8'h00;
    else count <= count + 1'b1;
  end
  initial $display("counter is ready");
endmodule
\end{CodeBlock}

\subsection*{目录树与特殊字符}
\begin{CodeBlock}[language=BookText]
chip
├── rtl  处理器设计
│   ├── core.sv
│   └── soc.sv
└── sim  仿真环境
ASCII: 0 O 1 l I  _ - --  ' "  $ # % & < >
\end{CodeBlock}

\clearpage
\section*{通用语言与硬件设计}
\subsection*{C：关键字、注释、字符串}
\begin{CodeBlock}[language=BookC]
#include <stdint.h>
// 中文注释应完整显示，并与英文注释保持同色。
uint32_t read_value(volatile uint32_t *address) {
    if (address == NULL) return 0;
    printf("value = %u\n", *address);
    return *address;
}
\end{CodeBlock}

\subsection*{Scala / Chisel / SpinalHDL}
\begin{CodeBlock}[language=BookScala]
val count = RegInit(0.U(8.W))
when (io.enable) {
  count := count + 1.U
}.otherwise {
  count := 0.U
}
\end{CodeBlock}

\subsection*{SystemVerilog}
\begin{CodeBlock}[language=BookSystemVerilog]
logic [31:0] data;
always_ff @(posedge clk) begin
  if (!reset_n) data <= '0;
  else data <= next_data;
end
\end{CodeBlock}

\subsection*{Shell}
\begin{CodeBlock}[language=BookShell]
# 编译命令保持原始内容，不由 LaTeX 执行。
export ARCH=loongarch
make -j$(nproc)
echo "Build completed"
\end{CodeBlock}

\clearpage
\section*{本书专用语言规则}
\subsection*{LoongArch 汇编}
\begin{CodeBlock}[language=BookLoongArch]
.text
.globl reset_entry
reset_entry:
    li.w    $t0, 0x1c000000  # 初始化地址
    ld.w    $a0, $t0, 0
    addi.w  $a0, $a0, 1
    st.w    $a0, $t0, 0
    ertn
\end{CodeBlock}

\subsection*{设备树 DTS}
\begin{CodeBlock}[language=BookDTS]
uart0: serial@1fe001e0 {
    compatible = "ns16550a";
    reg = <0x1fe001e0 0x8>;
    interrupts = <GIC_SPI 2 IRQ_TYPE_LEVEL_HIGH>;
    status = "okay";
};
\end{CodeBlock}

\subsection*{EDA Tcl}
\begin{CodeBlock}[language=BookTcl]
# 创建时钟并设置输入延迟。
set target_library "typical.db"
read_verilog design.v
create_clock -name clk -period 10 [get_ports clk]
set_input_delay -clock clk 3 [all_inputs]
compile -map_effort high
\end{CodeBlock}

\clearpage
\section*{配置、构建与补丁}
\subsection*{Kconfig / defconfig}
\begin{CodeBlock}[language=BookConfig]
config MEGA_SOC_FPGA
    bool "MEGA SoC for FPGA"
    default y
# CONFIG_DEBUG_INFO is not set
CONFIG_LOG_BUF_SHIFT=18
\end{CodeBlock}

\subsection*{Makefile}
\begin{CodeBlock}[language=BookMake]
CC = loongarch32r-linux-gnusf-gcc
all: kernel.o
kernel.o: kernel.c
	$(CC) -c kernel.c -o kernel.o
\end{CodeBlock}

\subsection*{链接脚本}
\begin{CodeBlock}[language=BookLinker]
OUTPUT_FORMAT(elf32-loongarch)
MEMORY {
  ram (rwxa) : ORIGIN = 0x80000000, LENGTH = 128M
}
INCLUDE "section.ld"
\end{CodeBlock}

\clearpage
\section*{本书新增的高亮语言}
\subsection*{C++ 仿真模型}
\begin{CodeBlock}[language=BookCpp]
// 模型调用示例，保留类型、作用域符号与字符串。
class Model final {
public:
  uint32_t value = 0;
  void step() { value += 1; }
};
int main() {
  Model model;
  model.step();
  std::cout << "value = " << model.value << std::endl;
}
\end{CodeBlock}

\subsection*{JSON 配置（支持原稿的注释写法）}
\begin{CodeBlock}[language=BookJSON]
{
  // 解码器配置片段；原稿采用带注释的 JSON 示例。
  "module_name": "core_decoder",
  "enabled": true,
  "width": 32
}
\end{CodeBlock}

\subsection*{C Shell / 签核脚本}
\begin{CodeBlock}[language=BookCShell]
#!/bin/csh -f
# 保留与 Bash 不同的变量赋值写法。
set top = design
set gds = /design_path/design.gds.gz
if ( -e $gds ) then
  echo "Input exists: $gds"
endif
\end{CodeBlock}
\end{document}
